% GENERATED FILE. Edit canonical lessons, then run npm run generate:books.
% canonical-source-hash: fnv1a64:7dc0e4db6f2c9de9
% canonical-lessons: PT-C25-amigo, PT-C25-familia, PT-C25-homem, PT-C25-mulher

\chapter{The People You Introduce}
\label{ch:pt-people-you-introduce}
\hlchaptermodality{25}

\begin{chapteropening}
\textbf{By the end of this chapter:} \emph{I can name a friend, a family, a man and a woman in Portuguese with the right gender, and say which Latin root each of the four Romance sisters picked for 'woman'.}

The chapter's last lesson closes the set with a one-table review of all four people-words' Latin roots, shows mulher/mujer against femme (fēmina) and donna (domina) as a three-way Romance split, and reaches back to Chapter 3's introduction practice and Chapter 15's tenho falado to build Tenho falado com a minha familia.
\end{chapteropening}

\section[o amigo, a amiga]{o amigo, a amiga — \textquotedblleft{}one who is loved\textquotedblright{}}

\label{lesson:PT-C25-amigo}



\begin{quote}
New chapter: the people you'd introduce, not the things you'd order. First up is the word for the person you say \emph{prazer} to.
\end{quote}

\subsection*{You'll want to know: The word}
\begin{quote}
\textbf{o amigo} (a male friend) / \textbf{a amiga} (a female friend) \textbf{O meu amigo} / \textbf{a minha amiga} — \textquotedblleft{}my friend.\textquotedblright{}
\end{quote}

Like \emph{obrigado/obrigada} in Chapter 1, the ending carries the gender: \textbf{-o} for a male friend, \textbf{-a} for a female one — but here, unlike \emph{obrigado}, the \emph{-o}/\emph{-a} agrees with the \textbf{friend}, not the speaker.

\begin{cousinweb}
\textbf{amigo} $\leftarrow$ Latin \emph{\textbf{amīcus}} (\textquotedblleft{}friend\textquotedblright{}), built directly on the verb \emph{\textbf{amāre}}, \textquotedblleft{}to love.\textquotedblright{} A friend, in Latin, is literally \textbf{\textquotedblleft{}one who is loved.\textquotedblright{}} English kept the same root in \textbf{amicable}, \textbf{amiable}, and \textbf{amity} — all learned words for the everyday Portuguese one.

That closes a loop this book opened back in Chapter 3. You met \textbf{prazer} there — \textquotedblleft{}pleasure,\textquotedblright{} from \emph{placēre}, \textquotedblleft{}to please\textquotedblright{} — as the word you say on \textbf{meeting} someone. Now you have the word for the person you'd still be saying it to once they're no longer a stranger: \emph{placēre} gets you through the door, \emph{amāre} is what's supposed to be on the other side of it.
\end{cousinweb}

\subsection*{Your turn}
Say these aloud:

\begin{itemize}
\raggedright
  \item \textquotedblleft{}o amigo, a amiga\textquotedblright{} — the \emph{-o}/\emph{-a} names the friend's gender
  \item \textquotedblleft{}o meu amigo\textquotedblright{} · \textquotedblleft{}a minha amiga\textquotedblright{}
  \item \textquotedblleft{}amigo — amicable — amāre, to love\textquotedblright{}
  \item the pair — \textquotedblleft{}prazer, from placēre, to please\textquotedblright{} · \textquotedblleft{}amigo, from amāre, to love\textquotedblright{}
\end{itemize}

\begin{tcolorbox}[breakable,colback=teal!4,colframe=teal!35!black,title={Before you move on}]
Say \textquotedblleft{}a male friend\textquotedblright{} and \textquotedblleft{}a female friend.\textquotedblright{} (\textbf{O amigo, a amiga}.) What Latin verb is \emph{amigo} built on, and what does it mean? (Latin \emph{\textbf{amāre}}, \textquotedblleft{}\textbf{to love}\textquotedblright{} — a friend is \textquotedblleft{}one who is loved.\textquotedblright{}) Name the two Chapter 24 foods that crossed from Asia through Arabic. (\textbf{Arroz, açúcar}.) Next: the word for everyone in that house together.
\end{tcolorbox}
\section[a família]{a família — a word that used to mean \textquotedblleft{}the servants\textquotedblright{}}

\label{lesson:PT-C25-familia}



\begin{quote}
A word that looks like the easiest cognate in the whole book — and whose Latin meaning would surprise you.
\end{quote}

\subsection*{You'll want to know: The word}
\begin{quote}
\textbf{a família} — the family. \textbf{Feminine.} \textbf{A minha família} — \textquotedblleft{}my family.\textquotedblright{}
\end{quote}

Spelled almost exactly like English \emph{family}, stressed almost the same way — one of the plainest true cognates this book teaches.

\begin{cousinweb}
\textbf{família} $\leftarrow$ Latin \emph{\textbf{familia}} — but Classical Latin \emph{familia} did not mean quite what you'd guess. It named the whole \textbf{household} under one head (the \emph{pater familias}), and that household's core members were the \emph{\textbf{famuli}} — its \textbf{servants and slaves}. Blood relatives were part of a \emph{familia} too, but the word's original centre of gravity was \textbf{who lived under one roof and one authority}, not \textbf{who you were related to}. Only over centuries did the sense narrow to the one you already know.

So when Chapter 10 taught you \textbf{o pai, a mãe} — \textquotedblleft{}father, mother\textquotedblright{} — those two words named the \emph{heads} of a \emph{familia} in the Roman sense: the people in charge of the household, servants included. And Chapter 10's \textbf{os pais} (\textquotedblleft{}the parents,\textquotedblright{} literally \textquotedblleft{}the fathers\textquotedblright{}) and \textbf{os irmãos} (\textquotedblleft{}brothers,\textquotedblright{} from \emph{germānus}, \textquotedblleft{}of the same stock\textquotedblright{}) are exactly the members a \emph{familia} still holds onto now that the servants have left the word behind — the same household whose table Chapter 24 stocked with \textbf{arroz} and \textbf{açúcar}.
\end{cousinweb}

\subsection*{Your turn}
Say these aloud:

\begin{itemize}
\raggedright
  \item \textquotedblleft{}a família, a minha família\textquotedblright{}
  \item \textquotedblleft{}família — a Roman household, servants included\textquotedblright{}
  \item the family words it holds — \textquotedblleft{}os pais, os irmãos\textquotedblright{} — Chapter 10
  \item \textquotedblleft{}o pai, a mãe\textquotedblright{} were the \emph{heads} of a \emph{familia}, not its whole definition
\end{itemize}

\begin{tcolorbox}[breakable,colback=teal!4,colframe=teal!35!black,title={Before you move on}]
Say \textquotedblleft{}the family\textquotedblright{} and \textquotedblleft{}my family.\textquotedblright{} (\textbf{A família, a minha família}.) What did Latin \emph{familia} actually name, and who were its core members? (The whole \textbf{household}, headed by one person — its \emph{famuli} were \textbf{servants}.) What does \emph{os pais} mean literally, and what does \emph{os irmãos} mean? (Literally \textquotedblleft{}\textbf{the fathers}\textquotedblright{}; \textquotedblleft{}\textbf{brothers},\textquotedblright{} from \emph{germānus}, Chapter 10.) Which Chapter 24 sweetener shares that same household table? (\textbf{Açúcar}.) Next: one of the two people who used to head that household.
\end{tcolorbox}
\section[o homem]{o homem — \textquotedblleft{}human being,\textquotedblright{} narrowed to half of one}

\label{lesson:PT-C25-homem}



\begin{quote}
A word that looks like it should mean \textquotedblleft{}human\textquotedblright{} in general — and in its Latin root, it did.
\end{quote}

\subsection*{You'll want to know: The word}
\begin{quote}
\textbf{o homem} — the man. \textbf{Masculine.} \textbf{O homem tem trinta anos.} — \textquotedblleft{}The man is thirty.\textquotedblright{} (Reusing \emph{ter} + \emph{anos} from Chapter 14 — never \emph{ser}.)
\end{quote}

The \textbf{h} is silent, as in every Portuguese word that has one. The ending \textbf{-em} is a nasal, the same nasal sound you already own from \emph{também} and \emph{tem}.

\begin{cousinweb}
\textbf{homem} $\leftarrow$ Latin \emph{\textbf{homō, hominis}} — and Classical Latin \emph{homō} meant \textbf{\textquotedblleft{}human being,\textquotedblright{}} with \textbf{no sex specified at all}. English kept that broad sense whole, in \textbf{human}, \textbf{humane}, and \textbf{hominid} — a \emph{hominid} is any member of the human family, not \textquotedblleft{}a member of the family who is male.\textquotedblright{}

Portuguese narrowed the word down to mean specifically \textbf{a man}, the way Spanish did with \emph{hombre} and French did with \emph{homme} — three sisters independently taking the same general Latin word and pointing it at one sex. So an English speaker's instinct on first meeting \emph{homem} is \textbf{half right}: the root really did mean \textquotedblleft{}human,\textquotedblright{} but the Portuguese word in front of you no longer does.
\end{cousinweb}

\subsection*{Your turn}
Say these aloud:

\begin{itemize}
\raggedright
  \item \textquotedblleft{}o homem\textquotedblright{} — silent \emph{h}, nasal \emph{-em}
  \item \textquotedblleft{}O homem tem trinta anos\textquotedblright{} — Chapter 14's \emph{ter} + \emph{anos}, not \emph{ser}
  \item \textquotedblleft{}homo, hominis — human being\textquotedblright{} — then \textquotedblleft{}homem — narrowed to man\textquotedblright{}
  \item the three sisters — \textquotedblleft{}Portuguese homem, Spanish hombre, French homme\textquotedblright{}
\end{itemize}

\begin{tcolorbox}[breakable,colback=teal!4,colframe=teal!35!black,title={Before you move on}]
Say \textquotedblleft{}the man\textquotedblright{} and give him an age using Chapter 14's pattern. (\textbf{O homem}; \textbf{o homem tem trinta anos} — \emph{ter}, never \emph{ser}.) What did Latin \emph{homō} actually mean, with what English word keeping that sense whole? (\textbf{\textquotedblleft{}Human being,\textquotedblright{}} sex unspecified — English \textbf{human}, \textbf{hominid}.) What happened to that sense in Portuguese, Spanish and French? (\textbf{Narrowed} to \textquotedblleft{}man\textquotedblright{} in all three — \emph{homem}, \emph{hombre}, \emph{homme}.) Next: the word Portuguese built from a completely different Latin root instead of narrowing this one.
\end{tcolorbox}
\section[a mulher]{a mulher — the fourth sister, a fourth root}

\label{lesson:PT-C25-mulher}



\begin{quote}
The last person-word of the chapter — and the payoff of the whole set. Portuguese, Spanish, French and Italian each say \textquotedblleft{}woman\textquotedblright{} with a \textbf{different} Latin word.
\end{quote}

\subsection*{You'll want to know: The word}
\begin{quote}
\textbf{a mulher} — the woman. \textbf{Feminine.} \textbf{A mulher tem trinta anos.} — \textquotedblleft{}The woman is thirty.\textquotedblright{}
\end{quote}

The \textbf{-lh-} is the palatal \emph{ly} sound you already own from \emph{trabalhar} (Ch. 5) and \emph{vinho} (Ch. 11).

\begin{cousinweb}
\textbf{mulher} $\leftarrow$ Latin \emph{\textbf{mulier, mulieris}} — twin of Spanish \textbf{mujer}, both straight descendants of the same word.

French and Italian did not follow. \textbf{French femme} comes from a \textbf{different} Latin word, \emph{\textbf{fēmina}} (the same root behind English \emph{feminine}, \emph{female}). \textbf{Italian donna} comes from a \textbf{third} Latin word again — \emph{\textbf{domina}}, \textquotedblleft{}\textbf{lady of the house, mistress}\textquotedblright{} (the same root as English \emph{dominate} and \emph{dame}). So the sisters split three ways on this one concept: Iberian Romance keeps \emph{mulier}, French kept \emph{fēmina}, Italian switched to \emph{domina} entirely.
\end{cousinweb}

\begin{grammarlens}[title={What you've built: The chapter in one table}]
\noindent
\begin{tabularx}{\linewidth}{@{}>{\raggedright\arraybackslash}X>{\raggedright\arraybackslash}X>{\raggedright\arraybackslash}X@{}}
\toprule
\textbf{Portuguese} & \textbf{English cousin} & \textbf{Latin root} \\
\midrule
\textbf{amigo/amiga} & amicable & \emph{amīcus}, $\leftarrow$ \emph{amāre} \textquotedblleft{}to love\textquotedblright{} \\
\textbf{família} & family & \emph{familia}, \textquotedblleft{}household\textquotedblright{} \\
\textbf{homem} & human, hominid & \emph{homō}, \textquotedblleft{}human being\textquotedblright{} \\
\textbf{mulher} & (none direct) & \emph{mulier} \\
\bottomrule
\end{tabularx}

Not one Romance sister names \textquotedblleft{}woman\textquotedblright{} the same way twice — the same \textquotedblleft{}not one of them agrees\textquotedblright{} shape Chapter 13 found in Portuguese's own colours. Hand Chapter 3's \emph{Prazer} to any of these four people (\textbf{Prazer, sou amigo da mulher}), and give Chapter 15's \textbf{tenho falado} — \textquotedblleft{}I've \textbf{been} speaking,\textquotedblright{} never \textquotedblleft{}I have spoken\textquotedblright{} — someone to be about: \emph{Tenho falado com a minha família.}
\end{grammarlens}

\subsection*{Your turn}
Say these aloud:

\begin{itemize}
\raggedright
  \item \textquotedblleft{}a mulher\textquotedblright{} — the \emph{-lh-} from \emph{trabalhar} and \emph{vinho}
  \item \textquotedblleft{}mulher — mujer\textquotedblright{} the Iberian twins
  \item \textquotedblleft{}femme (fēmina), donna (domina), mulher (mulier)\textquotedblright{} — three Latin roots, three sisters
  \item \textquotedblleft{}Tenho falado com a minha família\textquotedblright{} — Chapter 15's compound, now with someone to say it about
\end{itemize}

\begin{tcolorbox}[breakable,colback=teal!4,colframe=teal!35!black,title={Before you move on}]
Say \textquotedblleft{}the woman.\textquotedblright{} (\textbf{A mulher}.) What Latin word is it from, and what's its Spanish twin? (\emph{\textbf{Mulier}} — Spanish \textbf{mujer}.) What two different Latin words did French and Italian pick instead, and what do they mean? (French \textbf{femme} $\leftarrow$ \emph{fēmina}, \textquotedblleft{}female\textquotedblright{}; Italian \textbf{donna} $\leftarrow$ \emph{domina}, \textquotedblleft{}lady of the house.\textquotedblright{}) What does \emph{tenho falado com a minha família} mean, and why not \textquotedblleft{}I have spoken\textquotedblright{}? (\textbf{\textquotedblleft{}I've been talking with my family,\textquotedblright{}} repeatedly, since Chapter 15's compound never took over the plain past the way Chapter 3's introduction script eventually got a formal twin too.) Next chapter: parts of that same family's bodies.
\end{tcolorbox}
